\subsection{Factibilidad}
 
El estudio de factibilidad es un instrumento que sirve para orientar la toma de decisiones en la evaluación de un proyecto. Se formula con base en información que tiene la menor incertidumbre posible para medir las posibilidades de éxito o fracaso, apoyándose en él se tomará la decisión de proceder o no con su implementación.
 
% ─────────────────────────────────────────────
\subsubsection{Factibilidad Técnica}
% ─────────────────────────────────────────────
 
La factibilidad técnica estudia la posibilidad tecnológica para que el proyecto pueda llevarse a cabo satisfactoriamente con el menor riesgo posible. Las herramientas tecnológicas se dividen en software y hardware.
 
\paragraph{Software}
 
El software es una parte esencial en la elaboración del proyecto, ya que con este se desarrollará la aplicación web y el servidor. Todas las herramientas seleccionadas tienen costo cero para el equipo, incluyendo GitHub en su plan gratuito, el cual cubre repositorios privados ilimitados, colaboradores ilimitados y 2{,}000 minutos mensuales de CI/CD, suficientes para un equipo de tres personas.
 
\begin{table}[H]
\centering
\caption{Herramientas de software para el desarrollo del proyecto.}
\label{tab:software}
\renewcommand{\arraystretch}{1.3}
\begin{tabularx}{\textwidth}{|l|X|l|r|r|}
\hline
\textbf{Herramienta} & \textbf{Uso en el proyecto} & \textbf{Tipo} & \textbf{Costo/mes} & \textbf{11 meses} \\
\hline
React / Next.js   & Front-End: panel tutor, panel alumno, lienzo digital  & Libre & \$0 & N/A \\
\hline
FastAPI (Python)  & Back-End: endpoints REST, lógica de negocio, JWT      & Libre & \$0 & N/A \\
\hline
SQL Server Express & Base de datos: tablas, triggers, vistas, restricciones & Libre & \$0 & N/A \\
\hline
Docker            & Contenedorización y despliegue                        & Libre & \$0 & N/A \\
\hline
VS Code           & Editor de código fuente                               & Libre & \$0 & N/A \\
\hline
GitHub (Free)     & Control de versiones, repos privados, CI/CD           & Libre & \$0 & N/A \\
\hline
GCP / AWS         & Hosting: prueba piloto en tier gratuito               & Libre & \$0 & N/A \\
\hline
\multicolumn{3}{|r|}{\textbf{Total}} & \multicolumn{2}{r|}{\textbf{\$0 MXN}} \\
\hline
\end{tabularx}
\end{table}
 
\paragraph{Hardware}
 
El hardware es fundamental para el desarrollo del proyecto. A continuación se presentan las especificaciones de los equipos de cómputo de cada integrante del equipo, los cuales serán utilizados durante todo el desarrollo.
 
\begin{table}[H]
\centering
\caption{Características de los equipos de cómputo y dispositivos de prueba.}
\label{tab:hardware}
\renewcommand{\arraystretch}{1.3}
\begin{tabularx}{\textwidth}{|X|X|c|c|c|r|}
\hline
\textbf{Equipo} & \textbf{Procesador} & \textbf{RAM} & \textbf{Almac.} & \textbf{SO} & \textbf{Valor ref. (MXN)} \\
\hline
Lenovo IdeaPad  & AMD Ryzen 7 & 16 GB & 512 GB SSD & Arch Linux & \$18,000 \\
\hline
MacBook Air     & Apple M1 & 8 GB & 256 GB SSD & macOS & \$28,000 \\
\hline
HP Spectre x360 & Intel Core i7 & 16 GB & 512 GB SSD & Windows 11 & \$25,000 \\
\hline
iPad / Tablets  & N/A & N/A & N/A & iPadOS/Android & \$10,000 \\
\hline
\multicolumn{5}{|r|}{\textbf{Total (valor de referencia — ya en posesión del equipo)}} & \textbf{\$81,000} \\
\hline
\end{tabularx}
\end{table}
 
Dadas las características de las herramientas tanto en software como en hardware, se puede observar que el proyecto puede llevarse a cabo correctamente. El stack tecnológico es coherente y moderno: React/Next.js cubre el front-end educativo, FastAPI con Python concentra la lógica de inteligencia artificial (OCR y NLP), y SQL Server garantiza integridad referencial con soporte a triggers y vistas. El Ryzen 7 y el M1 son los equipos más capaces para ejecutar los modelos localmente durante el desarrollo.
 
% ─────────────────────────────────────────────
\subsubsection{Factibilidad Operativa}
% ─────────────────────────────────────────────
 
La factibilidad operativa comprende una determinación de la probabilidad para que un proyecto se realice o funcione como se supone, considerando si se cuenta con el personal capacitado y los recursos humanos necesarios. Al realizar el análisis se calcula tanto el tiempo de esfuerzo como las horas dedicadas al sistema.
 
\begin{table}[H]
\centering
\caption{Resumen de días hábiles para el desarrollo del proyecto.}
\label{tab:dias}
\renewcommand{\arraystretch}{1.3}
\begin{tabularx}{\textwidth}{|l|r|r|r|r|r|r|}
\hline
\textbf{Mes} & \textbf{Días} & \textbf{Fin de semana} & \textbf{Días no laborales} & \textbf{Días hábiles} & \textbf{Hrs/día} & \textbf{Hrs totales} \\
\hline
Enero      & 31 & 9 & 1 & 21 & 4 & 84  \\
\hline
Febrero    & 28 & 8 & 1 & 19 & 4 & 76  \\
\hline
Marzo      & 31 & 8 & 1 & 22 & 4 & 88  \\
\hline
Abril      & 30 & 8 & 2 & 20 & 4 & 80  \\
\hline
Mayo       & 31 & 9 & 1 & 21 & 4 & 84  \\
\hline
Junio      & 30 & 8 & 0 & 22 & 4 & 88  \\
\hline
Julio      & 31 & 9 & 0 & 22 & 4 & 88  \\
\hline
Agosto     & 31 & 9 & 0 & 22 & 4 & 88  \\
\hline
Septiembre & 30 & 8 & 1 & 21 & 4 & 84  \\
\hline
\textbf{Total} & & & & \textbf{190} & & \textbf{760} \\
\hline
\end{tabularx}
\end{table}
 
Es necesaria la definición de roles en el desarrollo del proyecto para llevar una organización eficiente donde cada integrante desempeñe sus actividades sin generar retrasos. Dado que el equipo de trabajo cuenta con tres personas, es necesario que los integrantes desempeñen más de un rol a lo largo del proyecto.
 
\begin{table}[H]
\centering
\caption{Roles de trabajo del proyecto.}
\label{tab:roles}
\renewcommand{\arraystretch}{1.3}
\begin{tabularx}{\textwidth}{|l|X|r|}
\hline
\textbf{Rol} & \textbf{Descripción} & \textbf{Personas necesarias} \\
\hline
Líder de proyecto       & Encargado de la coordinación y planeación del proyecto.                                          & 1 \\
\hline
Analista / Arquitecto   & Realiza el análisis e diseño de la arquitectura del sistema y sus módulos.                       & 1 \\
\hline
Desarrollador Front-End & Desarrolla la interfaz de usuario con React/Next.js: paneles de tutor y alumno, lienzo digital.  & 1 \\
\hline
Desarrollador Back-End  & Implementa los endpoints REST con FastAPI, lógica de negocio, autenticación JWT y BCrypt.        & 1 \\
\hline
Especialista IA         & Integra y ajusta los módulos de OCR y NLP para análisis caligráfico y corrección ortográfica.    & 1 \\
\hline
Tester / QA             & Realiza las pruebas funcionales y de integración para validar el correcto funcionamiento.        & 1 \\
\hline
\end{tabularx}
\end{table}
 
A continuación se muestra la distribución del porcentaje de participación para cada rol y el promedio de días laborales correspondiente, tomando como base 190 días hábiles de desarrollo.
 
\begin{table}[H]
\centering
\caption{Porcentaje de participación de cada rol del proyecto.}
\label{tab:participacion}
\renewcommand{\arraystretch}{1.3}
\begin{tabularx}{\textwidth}{|X|r|r|}
\hline
\textbf{Rol} & \textbf{Porcentaje de participación} & \textbf{Días laborales} \\
\hline
Líder de proyecto       & 15\% & 28.5 días \\
\hline
Analista / Arquitecto   & 15\% & 28.5 días \\
\hline
Desarrollador Front-End & 20\% & 38.0 días \\
\hline
Desarrollador Back-End  & 20\% & 38.0 días \\
\hline
Especialista IA (OCR + NLP) & 20\% & 38.0 días \\
\hline
Tester / QA             & 10\% & 19.0 días \\
\hline
\end{tabularx}
\end{table}
 
La factibilidad operativa del proyecto se muestra positiva. Los días totales calculados son una estimación considerando una dedicación de 4 horas diarias promedio por integrante. Consideramos que el tiempo para el desarrollo tiene suficiente holgura siempre y cuando se sigan las actividades y los calendarios acordados, y se priorice la integración temprana entre módulos para detectar problemas antes de la prueba piloto con 25 alumnos.
 
% ─────────────────────────────────────────────
\subsubsection{Factibilidad Económica}
% ─────────────────────────────────────────────
 
Se toma en consideración el capital disponible para invertir en el desarrollo del proyecto. Se determina el presupuesto de costos de los recursos técnicos, humanos y materiales tanto para el desarrollo como para la implantación. El sueldo de referencia utilizado es de \$20{,}000 MXN mensuales por desarrollador de nivel junior a intermedio en México para el año 2025 \cite{occmundial2025salarios}.
 
En primer lugar, se presentan los costos de las herramientas de software, los cuales son en su totalidad de costo cero para el equipo de desarrollo.
 
\begin{table}[H]
\centering
\caption{Estimación de costos de herramientas de software.}
\label{tab:costos-software}
\renewcommand{\arraystretch}{1.3}
\begin{tabularx}{\textwidth}{|X|r|r|}
\hline
\textbf{Herramienta} & \textbf{Costo mensual (MXN)} & \textbf{Costo por 9 meses (MXN)} \\
\hline
React / Next.js, FastAPI, SQL Server Express, Docker, VS Code, Figma, Postman, GitHub Free, GCP/AWS (tier gratuito) & \$0 & \$0 \\
\hline
\textbf{Total} & \textbf{\$0} & \textbf{\$0} \\
\hline
\end{tabularx}
\end{table}
 
A continuación se presentan los costos de los servicios necesarios para llevar a cabo el proyecto.
 
\begin{table}[H]
\centering
\caption{Costos de servicios.}
\label{tab:servicios}
\renewcommand{\arraystretch}{1.3}
\begin{tabularx}{\textwidth}{|X|r|r|}
\hline
\textbf{Servicio} & \textbf{Costo por mes (MXN)} & \textbf{Costo por 9 meses (MXN)} \\
\hline
Electricidad e internet (3 integrantes × \$800) & \$2{,}400 & \$21{,}600 \\
\hline
Infraestructura en nube (GCP / AWS — tier gratuito en piloto) & \$0 & \$0 \\
\hline
\textbf{Total} & \textbf{\$2{,}400} & \textbf{\$21{,}600} \\
\hline
\end{tabularx}
\end{table}
 
Finalmente, en la Tabla \ref{tab:costos-total} se presenta el resumen de costos totales del proyecto.
 
\begin{table}[H]
\centering
\caption{Resumen de costos totales del proyecto.}
\label{tab:costos-total}
\renewcommand{\arraystretch}{1.3}
\begin{tabularx}{\textwidth}{|X|r|}
\hline
\textbf{Concepto} & \textbf{Costo total (MXN)} \\
\hline
Sueldos — costo comercial equivalente (3 personas × 9 meses × \$20{,}000) & \$540{,}000 \\
\hline
Equipo de cómputo (valor de referencia — ya en posesión del equipo)        & \$61{,}000 \\
\hline
Servicios (electricidad e internet, 9 meses)                               & \$21{,}600 \\
\hline
Herramientas de software                                                   & \$0        \\
\hline
Infraestructura en nube                                                    & \$0        \\
\hline
\textbf{Total}                                                             & \textbf{\$622{,}600} \\
\hline
\end{tabularx}
\end{table}
 
Con base en la tabla anterior, se puede concluir que el costo comercial equivalente del proyecto asciende a \$622{,}600 MXN. Sin embargo, al ser este un trabajo terminal, los sueldos (que representan el 86.7\% del total) son absorbidos por el propio equipo como parte de su formación académica, y el hardware ya se encuentra en posesión de los integrantes. El desembolso directo real del equipo se limita a los servicios básicos de electricidad e internet, estimados en \$21{,}600 MXN durante los 9 meses de desarrollo, siendo el resto de las herramientas completamente gratuitas.